\documentclass{article}

\textwidth=6.5in
\textheight=8.5in
\voffset=-54pt
\oddsidemargin=0in
\evensidemargin=0in
\setlength{\parskip}{8pt}
\setlength{\parindent}{0pt}

\usepackage{workbook}

\usepackage[]{handout}

\begin{document}

\begin{center}
  \large\textsc{Problem Set 30 - Integration Strategy}
  
      \pspace{12pt}
  \begin{problemonly}\large Name: \rule{5cm}{0.15mm}\\ \end{problemonly}
\end{center}

\begin{grading}
    \begin{enumerate}
        \item 17 points 2/2/5/3/5
        \item 10 points, 2/3/3/2
        \item 25 points, 5 for each part
    \end{enumerate}
    Total: 52 points
\end{grading}

\begin{enumerate}
    \item 
    Find or evaluate the following integrals.
    \begin{enumerate}
        \item 
        \begin{problem}[\vfill]
            $\displaystyle\int \frac{1}{2x+1}\, dx$
        \end{problem}

        \begin{solution}
            We can use the substitution $u=2x+1$, $du = 2\,dx$.
            \begin{align*}
                \int \frac{1}{2x+1}\, dx &= \frac{1}{2} \int \frac{1}{u} \, du \\
                &= \frac{1}{2} \ln|u| +C \\
                &= \frac{1}{2}\ln|2x+1| +C 
            \end{align*}
        \end{solution}

        \item 
        \begin{problem}[\vfill]
            $\displaystyle\int_2^6 \left|\tfrac{1}{2}t-2\right|\, dt$
        \end{problem}

        \begin{solution}
            We can evaluate this definite integral using geometry.
            \begin{center}
            \begin{tikzpicture}
              \begin{axis}[cycle list=black, axis equal=true, xtick={2, 4, 6},
                ytick={1}, width=2.5in]
                \addplot+[domain=2:6, plusarea] {abs(0.5*x-2)} \closedcycle;
                \addplot+[domain=0:6.5] {abs(0.5*x-2)};
              \end{axis}
            \end{tikzpicture} 
            \end{center}
            The signed area under the graph of the function is 
            $\displaystyle\int_2^6 \left|\tfrac{1}{2}t-2\right|\, dt=2.$
        \end{solution}

    

     \item 
        \begin{problem}[\vfill\newpage]
            $\displaystyle\int \frac{x+1}{x^2+1}\, dx$
        \end{problem}

        \begin{solution}
            We can rewrite this integral as 
            \[\int \frac{x+1}{x^2+1}\, dx = \int \frac{x}{x^2+1}\,dx + 
            \int \frac{1}{x^2+1}\,dx.\]
            The first integral can be done using the substitution $u=x^2+1$, 
            $du = 2x\,dx$. The second is an elementary antiderivative.
            \begin{align*}
            \int \frac{x+1}{x^2+1}\, dx &= \frac{1}{2}\int \frac{1}{u}\,du + 
            \int \frac{1}{x^2+1}\,dx \\
            &= \tfrac{1}{2}\ln|u| + \arctan(x) + C \\
            &= \tfrac{1}{2}\ln(x^2+1) + \arctan(x) +C.
            \end{align*}
        \end{solution}

        \item 
        \begin{problem}[\vfill]
            $\displaystyle\int_1^2 \frac{3x}{x^2+8}\, dx$
        \end{problem}

      \begin{solution} % originally by Henry Shin, modified by Janet
        Let $\displaystyle u = x^2 + 8$; then, $du = 2x~dx$.  When $x = 1$,
        $u = 1^2 + 8 = 9$; when $x = 2$, $u = 2^2 + 8 = 12$.  So, 
        \begin{align*}
          \int_1^2 \frac{3x}{x^2 + 8}~dx & =  
          \int_1^2 \frac{3}{2} \cdot \frac{2x}{x^2 + 8}~dx \\
          & = \int_9^{12} \frac{3}{2} \cdot \frac{1}{u}~du \\
          & = \left. \frac{3}{2} \ln |u| \right|_9^{12} \\
          & = \boxed{\frac{3}{2} (\ln 12 - \ln 9)} \\
          &= \frac{3}{2} \ln \frac{4}{3} 
        \end{align*}
      \end{solution}
    
    \item 
        \begin{problem}[\vfill]
            $\displaystyle \int \ln(3x+2)\, dx$
        \end{problem}
        
        \begin{solution}
            We can first do the substitution $w = 3x+2$, $dw = 3\,dx$,
            giving us
            \[\int \ln(3x+2) \,dx = \frac{1}{3}\int \ln(w)\,dw.\]
            Then we can do integration by parts with
            \begin{align*}
                u = \ln(w) \qquad & dv = dw \\
                du = \frac{1}{w}\,dw \qquad & v = w
            \end{align*}
            giving us
            \begin{align*}
            \int \ln(3x+2)\, dx &= \frac{1}{3}\int \ln(w)\,dw \\
            &= \frac{1}{3}\left[w\ln(w) - \int 1\,dw\right] \\
            &= \tfrac{1}{3}w\ln(w) - \tfrac{1}{3}w +C \\
            &= \tfrac{1}{3}(3x+2)\ln(3x+2) - \tfrac{1}{3}(3x+2) +C
            \end{align*}
        \end{solution}

        \item 
        \begin{problem}[\vfill]
            $\displaystyle \int \sin\sqrt{x}\, dx$
        \end{problem}
        
        \begin{solution}
            It's not immediately obvious what method we should use here, but something we might try is to make the substitution $u=\sqrt{x}$, $du=\frac{1}{2\sqrt{x}}\,dx$. To change our integral to be completely in terms of $u$, $\sin\sqrt{x}$ will turn into $\sin u$, and since $x=u^2$, $dx=2u\,du$. So we have
            \begin{align*}
                \int \sin\sqrt{x} \, dx &= 2\int u \sin u \,du
                \intertext{The integral on the right-hand side can be evaluated using integration by parts,}
                &= 2\Big({-u}\cos u + \int \cos u \, du\Big) \\
                &= 2({-u}\cos u + \sin u) + C
                \intertext{Changing back to $x$:}
                &= 2(-\sqrt{x}\cos\sqrt{x} + \sin\sqrt{x}) + C
            \end{align*}
        \end{solution}

    
\pspace{\newpage}

        \item 
        \begin{problem}[\vfill]
            $\displaystyle \int_0^2 \sqrt{4-x}\, dx$
        \end{problem}
        
        \begin{solution}
            We can use a substitution to turn this into an integral we can evaluate using the power rule, $u=4-x$, $du=-dx$. 
            \begin{align*}
                \int_0^2 \sqrt{4-x}\, dx &= -\int_4^2 \sqrt{u}\,du \\
                &= \int_2^4 \sqrt{u}\,du \\
                &= \big[\tfrac{2}{3}u^{3/2}\big]_2^4 \\
                &= \tfrac{2}{3}(4^{3/2}-2^{3/2}) \\
                &= \tfrac{2}{3}(8-2\sqrt{2}) \\
                &= \tfrac{4}{3}(4-\sqrt{2})
            \end{align*}
        \end{solution}

        \item 
        \begin{problem}[\vfill]
            $\displaystyle \int_0^2 \sqrt{4-x^2}\, dx$
        \end{problem}
        \begin{solution}
        $y=\sqrt{4-x^2}$ is the graph of the top half of a circle of radius $2$ centered at the origin. So using geometry, \[\int_0^2 \sqrt{4-x^2}\, dx=\pi\]
        \end{solution}
        
        \item 
        \begin{problem}[\vfill]
            $\displaystyle \int_0^2 x\sqrt{4-x^2}\, dx$
        \end{problem}
        \begin{solution}
            We can use a substitution to undo the Chain Rule. Let $u=4-x^2$, so $du=-2x\,dx$.
            \begin{align*}
                \int_0^2 x\sqrt{4-x^2}\,dx &= 
                -\tfrac{1}{2}\!\int_4^0\sqrt{u}\,du \\
                &= \tfrac{1}{2}\!\int_0^4 \sqrt{u}\, du \\
                &= \tfrac{1}{2}\big[\tfrac{2}{3}u^{3/2}\big]_0^4 \\
                &= \tfrac{1}{3}(4^{3/2}) \\
                &= \tfrac{8}{3}
            \end{align*}
        \end{solution}

    \end{enumerate}

    \pspace{\newpage}
\item 
\begin{grading}
10 points. 2/3/3/2
\end{grading}
 \emph{Preview of next class.}

Shown below are the graphs of $y=3x-x^2$ and $y=\tfrac{1}{2}x^2$, and the region $\mathcal{R}$ that the two curves enclose. We say that $\mathcal{R}$ is \emph{bounded} by $y=3x-x^2$ and $y=\tfrac{1}{2}x^2$. We are going to find the area of $\mathcal{R}$.

\begin{center}
\begin{tikzpicture}
\begin{axis}[xlabel=$x$, ylabel=$y$, xtick=\empty, ytick=\empty, clip=false, width=3in]
    \addplot[domain=-0.25:2.5] {3*x-x^2};
    \node[anchor=west] at (2.5,3.25) {$y=\tfrac{1}{2}x^2$};
    \addplot[domain=-0.5:2.5] {0.5*x^2};
    \node[anchor=west] at (2.5,1.25) {$y=3x-x^2$};
    \addplot[name path=A, domain=0:2] {3*x-x^2};
    \addplot[name path=B, domain=0:2] {0.5*x^2};
    \addplot[fill=gray, fill opacity=0.3] fill between[of=A and B];
    \node at (1,1.25) {\normalsize$\mathcal{R}$};
    \solonly{
    \addplot[closed] coordinates {(0,0) (2,2)};
    \node[below right] at (0,0) {$(0, 0)$};
    \node[xshift=16pt] at (2, 2) {$(2,2)$};
    }
\end{axis}
\end{tikzpicture}
\end{center}

\begin{enumerate}
    \item 
    \begin{problem}
        Find and label the points of intersection.
    \end{problem}
    \pspace{\vfill}

    \begin{solution}
        To find the points of intersection, we need to find the $x$-values for which the $y$-values of the two curves are the same. That is, we need to solve the following equation.
        \begin{align*}
            3x-x^2 &= \tfrac{1}{2}x^2 \\
            3x-\tfrac{3}{2}x^2 &= 0 \\
            -\tfrac{3}{2}x(x-2)&= 0 \\
            x&= 0, 2
        \end{align*}
        The $x$-coordinates of the points of intersection are $0$ and $2$, and we can plug these into the equation of either curve to find that the points of intersection are $(0, 0)$ and $(2, 2)$.
    \end{solution}

    \item 
    \begin{problem}
        Based on the diagram below, express the area of $\mathcal{R}$ as the \textbf{difference} of two definite integrals. Evaluate the definite integrals and give an exact value for the area of $\mathcal{R}$. 
    \end{problem}


\begin{tikzpicture}
\begin{axis}[xlabel=$x$, ylabel=$y$, xtick=\empty, ytick=\empty, clip=false, width=3in]
    \addplot[domain=-0.25:2.5] {3*x-x^2};
    \node[anchor=west] at (2.5,3.25) {$y=\tfrac{1}{2}x^2$};
    \addplot[domain=-0.5:2.5] {0.5*x^2};
    \node[anchor=west] at (2.5,1.25) {$y=3x-x^2$};
    \addplot[name path=A, domain=0:2] {3*x-x^2};
    \addplot[name path=B, domain=0:2] {0.5*x^2};
    \addplot[fill=gray, fill opacity=0.3] fill between[of=A and B];
    \node at (1,1.25) {\normalsize$\mathcal{R}$};
    \addplot[name path=X, draw=none, domain=0:2] {0};
    \addplot[fill=white, pattern=north west lines, fill opacity=0.3] fill between[of=B and X];
    \draw[dashed] (2,0) -- (2,2);
\end{axis}
\end{tikzpicture}
\pspace{\stretch{0.8}}

    \begin{solution}
    Based on the diagram, the area of $\mathcal{R}$ is
    \begin{align*}
        \int_0^2 (3x-x^2)\,dx - \int_0^2 \tfrac{1}{2}x^2\,dx
        &= \big[\tfrac{3}{2}x^2-\tfrac{1}{3}x^3\big]_0^2 - \big[\tfrac{1}{6}x^3\big]_0^2 \\
        &= (6-\tfrac{8}{3}) - (\tfrac{4}{3}) \\
        &= 2
    \end{align*}
    \end{solution}

\end{enumerate}

\pspace{\newpage}

Now we are going to approach finding the area of $\mathcal{R}$ a different way, by constructing an integral from the limit of a Riemann sum. The diagram below shows how we can approximate the area of $\mathcal{R}$ using a right Riemann sum with 4 equal sub-intervals, $R_4$.

  \begin{tikzpicture}
    \begin{groupplot}[group style={group size=2 by 1},xlabel=$x$, ylabel=$y$, xtick=\empty, ytick=\empty, clip=false, width=3in]
    \nextgroupplot
    \addplot[domain=-0.25:2.5] {3*x-x^2};

    \addplot[domain=-0.5:2.5] {0.5*x^2};
    \addplot[name path=A, domain=0:2] {3*x-x^2};
    \addplot[name path=B, domain=0:2] {0.5*x^2};
    \addplot[fill=gray, fill opacity=0.3] fill between[of=A and B];
    \draw[black, thick] (0.5, {0.5*0.5^2}) -- (0.5, {3*0.5 - 0.5*0.5});
        \draw[black, thick] (1, {0.5*1^2}) rectangle (1, {3*1 - 1*1});
        \draw[black, thick] (1.5, {0.5*1.5^2}) rectangle (1.5, {3*1.5 - 1.5*1.5});
    \nextgroupplot
    \addplot[domain=-0.25:2.5] {3*x-x^2};
    \node[anchor=west] at (2.5,3.25) {$y=\tfrac{1}{2}x^2$};
    \addplot[domain=-0.5:2.5] {0.5*x^2};
    \node[anchor=west] at (2.5,1.25) {$y=3x-x^2$};
    \addplot[name path=A, domain=0:2] {3*x-x^2};
    \addplot[name path=B, domain=0:2] {0.5*x^2};
    \fill[gray, fill opacity=0.3, draw=black, thick] (0, {0.5*0.5^2}) rectangle (0.5, {3*0.5 - 0.5*0.5});
        \fill[gray, fill opacity=0.3, draw=black,  thick] (0.5, {0.5*1^2}) rectangle (1, {3*1 - 1*1});
        \fill[gray, fill opacity=0.3, draw=black, thick] (1, {0.5*1.5^2}) rectangle (1.5, {3*1.5 - 1.5*1.5});
        \fill[gray, fill opacity=0.3, draw=black, thick] (1.5, {0.5*2^2}) rectangle (2, {3*2 - 2*2});
    \end{groupplot}
  \end{tikzpicture}

\begin{enumerate}[resume]
\item 
\begin{problem}
    Write down $R_4$ as a sum of four terms and using sigma notation. Evaluate the sum and check that it makes sense when compared to your answer to part (b).
\end{problem}
\pspace{\vfill}

\begin{solution}
If we partition the interval $[0,2]$ into 4 equal subintervals, then the width of each sub-interval will be $\Delta x = \tfrac{1}{2}$. To find $R_4$, we need the heights of the rectangles. 

From the diagram, we can see that the height of a rectangle is the \emph{difference} between the $y$-coordinate of the top curve and the $y$-coordinate of the bottom curve at the same $x$-coordinate. For example, the height of the first rectangle is $[3(\frac{1}{2})-(\frac{1}{2})^2] - [\frac{1}{2}(\frac{1}{2})^2]=\frac{9}{8}$. Since the equation of the top curve is $y=3x-x^2$ and the equation of the bottom curve is $y=\frac{1}{2}x^2$, the difference $3x-\frac{3}{2}x^2$ gives us the heights of the rectangles at some $x$-value.
% \begin{itemize}
% \item
% The height of the first rectangle is $[3(\frac{1}{2})-(\frac{1}{2})^2] - [\frac{1}{2}(\frac{1}{2})^2]=$.
% \end{itemize}
\begin{align*}
    R_4 &= \tfrac{1}{2}[3(\tfrac{1}{2})-\tfrac{3}{2}(\tfrac{1}{2})^2] + \tfrac{1}{2}[3(1)-\tfrac{3}{2}(1)^2] + 
    \tfrac{1}{2}[3(\tfrac{3}{2})-\tfrac{3}{2}(\tfrac{3}{2})^2]+
    \tfrac{1}{2}[3(2)-\tfrac{3}{2}(2)^2]
    \intertext{In sigma notation,}
    R_4 &= \sum_{k=1}^4 \tfrac{1}{2}\big[3(\tfrac{k}{2})-\tfrac{3}{2}(\tfrac{k}{2})^2\big] \\
    \intertext{Evaluating,}
    &= \tfrac{1}{2}(\tfrac{9}{8}) + \tfrac{1}{2}(\tfrac{3}{2}) + \tfrac{1}{2}(\tfrac{9}{8})+\tfrac{1}{2}(0) \\
    &= \tfrac{15}{8}
\end{align*}
This value makes sense, because we found in part (b) that the exact area is $2$, and $\frac{15}{8}\approx 2$.
\end{solution}

\item 
\begin{problem}
    What definite integral would $\lim\limits_{n\to\infty} R_n$ be? Does this match what you said in part (b)?
\end{problem}
\pspace{\stretch{0.3}}

\begin{solution}
    As the number of rectangles goes to infinity, we'll be integrating from $x=0$ to $x=2$, and the height of the rectangle at any $x$-value is $3x-\tfrac{3}{2}x^2$, so we have the integral 
    \[\int_0^2 (3x-\tfrac{3}{2}x^2)\,dx\]
    This makes sense compared to the answer to part (b) because the expressions are equivalent, using our properties of integrals.
    \[\int_0^2 (3x-x^2)\,dx - \int_0^2 \tfrac{1}{2}x^2\,dx = \int_0^2 (3x-\tfrac{3}{2}x^2)\,dx\]
\end{solution}
\end{enumerate}


% Suppose we'd like to calculate the area of the region 
%   $\displaystyle \int_3^8 f(x)\,dx$ using the definition of the integral
%   as a limit of right-hand sums.  We'll start by trying to write a formula
%   for $R_n$, the right-hand sum with $n$ rectangles.

%   \begin{tikzpicture}[declare function={f(\x)=2*sin(deg(\x+1))+\x/3+4;}] 
%     \begin{groupplot}[group style={group size=2 by 1}]
%     \nextgroupplot[ymin=0, hide y axis, xtick={3, 8}, width=3in, height=1.8in,
%       cycle list=black, clip=false]
%       \addplot+[domain=3:8, fill=white!95!black] {f(x)} \closedcycle; %
%       \addplot+[domain=2:9] {f(x)} node[right] {$f$}; %
%       \nextgroupplot[ymin=0, hide y axis, xtick={3, 3.625, ..., 8.01},
%         xticklabels={$x_0$, $x_1$, $x_2$, $\cdots$, , , , , $x_n$}, 
%         width=3in, height=1.8in, disabledatascaling, cycle list=black,
%         clip=false] 
%         \addplot+[domain=3:8, fill=white!95!black] {f(x)} \closedcycle; %
%         \addplot+[domain=2:9] {f(x)} node[right]{$f$}; %
%         \addplot+[vertslices] coordinates { (3, {f(3)}) (29/8,
%           {f(29/8)}) (17/4, {f(17/4)}) (39/8, {f(39/8)}) (11/2, {f(11/2)})
%           (49/8, {f(49/8)}) (27/4, {f(27/4)}) (59/8, {f(59/8)}) (8,
%           {f(8)}) }; %
%         \begin{scope}[font=\footnotesize, inner sep=0pt]
%           \draw (3, 0) node[yshift=-13pt, rotate=90] {$=$}; %
%           \draw (8, 0) node[yshift=-13pt, rotate=90] {$=$}; %
%           \draw (3, 0) node[below, yshift=-18pt] {$3$}; %
%           \draw (8, 0) node[below, yshift=-18pt] {$8$}; %
%         \end{scope}
%     \end{groupplot}
%   \end{tikzpicture}

%   \begin{enumerate}[topsep=0pt]
%   \item The first step is to \textbf{slice} the region into $n$ pieces of
%     equal thickness $\Delta x$, and we label $x_0 = 3, x_1, x_2, \ldots,
%     x_n = 8$ as usual.

%     \begin{enumerate}[topsep=0pt]
%     \item
%       \begin{problem}[\vfill]
%         Express $\Delta x$ in terms of $n$.
%       \end{problem}

%       \begin{solution}
%         $\Delta x = \dfrac{8 - 3}{n}$.
%       \end{solution}

%     \item
%       \begin{problem}[\vfill]
%         Express $x_k$ in terms of $\Delta x$.  (If you're puzzled, first
%         think about how you would express $x_1$ as $x_0 + \Delta x$. 
%         What about $x_2$? $x_3$? Then consider how you would express
%         $x_{37}$,
%         assuming there are at least $37$ slices.)
%       \end{problem}

%       \begin{solution}
%         $x_k = 3 + k \Delta x$.  (For example, $x_1 = 3 + \Delta x$ and
%         $x_2 = 3 + 2 \Delta x$.) 
%       \end{solution}

%     \end{enumerate}

%     \pspace{\newpage}

%   \item
%     \begin{problem}[\vfill]
%       We then \textbf{approximate} each slice as a rectangle.  Remember
%       that we had decided to do a right-hand sum, so the rectangles look
%       like this. On each interval, we represent the height on that interval
%       as approximately the height on the right hand endpoint of the interval.

%       \begin{tikzpicture}[declare function={f(\x)=2*sin(deg(\x+1))+\x/3+4;}] 
%         \begin{axis}[ymin=0, hide y axis, xtick={3, 3.625, ..., 8.01},
%           xticklabels={$x_0$, $x_1$, $x_2$, $\cdots$, , , , , $x_n$}, 
%           width=3in, height=1.8in, disabledatascaling, cycle list=black,
%           clip=false] 
%           \addplot+[domain=3:8, fill=white!95!black] {f(x)} \closedcycle; %
%           \addplot+[domain=2:9] {f(x)} node[right]{$f$}; %
%         \addplot+[vertslices] coordinates { (3, {f(3)}) (29/8,
%           {f(29/8)}) (17/4, {f(17/4)}) (39/8, {f(39/8)}) (11/2, {f(11/2)})
%           (49/8, {f(49/8)}) (27/4, {f(27/4)}) (59/8, {f(59/8)}) (8,
%           {f(8)}) }; %
%           \begin{scope}[font=\footnotesize, inner sep=0pt]
%             \draw (3, 0) node[yshift=-13pt, rotate=90] {$=$}; %
%             \draw (8, 0) node[yshift=-13pt, rotate=90] {$=$}; %
%             \draw (3, 0) node[below, yshift=-18pt] {$3$}; %
%             \draw (8, 0) node[below, yshift=-18pt] {$8$}; %
%           \end{scope}
%           \addplot+[ybar interval, fill=cyan, fill opacity=0.25,
%           draw=blue] coordinates { (3, {f(29/8)}) (29/8, {f(17/4)}) (17/4,
%             {f(39/8)}) (39/8, {f(11/2)}) (11/2, {f(49/8)}) (49/8,
%             {f(27/4)}) (27/4, {f(59/8)}) (59/8, {f(8)}) (8, {f(8)}) };
%         \end{axis}
%       \end{tikzpicture}

%       What is the area of the 1st rectangle? The 2nd rectangle?
%       What is the area of the $k$-th rectangle?  Express your answer in
%       terms of $\Delta x$ and $x_k$ or $x_{k-1}$ as appropriate.  (If you're
%       puzzled, first think about how you would express the area of the
%       $37$th rectangle, assuming there are at least $37$ slices.)
%     \end{problem}

%     \begin{solution}
%       The area of the $k$-th rectangle is $f(x_k) \Delta x$.  (For
%       example, the area of the first rectangle is $f(x_1) \Delta x$, and
%       the area of the second rectangle is $f(x_2) \Delta x$.)
%     \end{solution}

%   \item
%     \begin{problem}[\vfill]
%       Finally, we \textbf{sum} the areas of all the rectangles to get
%       $R_n$.  Write a formula for $R_n$; express your answer once in
%       $\sum$ notation and once without $\sum$ notation.
%     \end{problem}

%     \begin{solution}
%       The area of the first rectangle is $f(x_1) \Delta x$, the area of
%       the second rectangle is $f(x_2) \Delta x$, and so on.  So, $R_n =
%       f(x_1) \Delta x + f(x_2) \Delta x + f(x_3) \Delta x + \cdots +
%       f(x_n) \Delta x = \displaystyle \sum_{k=1}^n f(x_k) \Delta x$. 
%     \end{solution}

%   \item
%     \begin{problem}[\vfill]
%       Finally, we \textbf{take a limit} to get the exact value of
%       $\displaystyle \int_3^8 f(x)\,dx$.  What limit do we take?  (That
%       is, we let which variable approach what?)
%     \end{problem}

%     \begin{solution}
%       We want to take the limit as $n \to \infty$. 
%     \end{solution}

%   \end{enumerate}


% \item 
% \begin{grading}
%     Students are asked to reflect on each of the 5 problems listed below. Each response is worth 5 points, and for full credit to be awarded, they need to answer each of the questions that are asked in the directions.
% \end{grading}

% Since the start of the integration unit, we have built an understanding of the integral from multiple perspectives: we have reasoned about the integral by looking at a graph, investigating a rate function, and building a Riemann sum. 
% In this problem, you will look back and summarize the key points of the math featured in five of the problems we’ve worked on. In that process, we hope you make new connections between your observations about the idea of the integral so that it is a long-lasting tool that you can use in your future courses.

% For each of the problems listed below, review your work for this problem on either the problem set or worksheet and write 3-6 sentences discussing the following:
% What information was communicated and asked for in the problem (this is given and unknown information)? What techniques did you need to solve the problem? Why were they effective methods? What do you think the problem’s purpose is? That is, why do you think we included it in our coverage of the integral? What did it help you better understand about integration?

% The problems to reflect on are:
% \begin{multicols}{2}
% \begin{itemize}
%     \item Worksheet 21 \#1
%     \item Worksheet 21 \#4
%     \item Worksheet 24 \#2
%     \item Worksheet 24 \#4
%     \item PSet 24 \#4 
% \end{itemize}
% \end{multicols}




% Example: 
% Worksheet 23 \#4

% We were asked to calculate the definite integral of $x^2$ from 0 to 1. This was challenging because we can’t break down the integral into basic geometric shapes. We used the limit definition of the derivative, and evaluating this limit was a complicated process that involved sigma notation, some mathematical fancy facts about the sum of the first n square numbers, and L’Hopital’s Rule. At the end, we decided we needed another way to calculate definite integrals. Looking back, this type of problem is very limited because you can’t complete it without the mathematical fact. I can see that the Fundamental Theorem of Calculus addresses this problem; it is much easier to calculate the integral using FTOC. 


  
% \end{enumerate}

% \newpage

% \textit{(This page is blank for you to write your responses to Problem 3)}

% \newpage 

% \textit{(This page is blank for you to write your responses to Problem 3)}

% \newpage 

% \;

\end{enumerate}
\psreflection{
Optional reading for next class is \S 27.2 in Gottlieb.
}
\end{document}